\chapter{The Ramsey Model}\label{ch:b2:ramsey-model}

In the Solow model of the Year 1 volume the saving rate is a fixed number. In the data it moves, and sometimes a great
deal. The chart shows gross domestic saving as a share of GDP in three Asian economies over six decades.

\begin{openingchart}{Gross domestic saving as a share of GDP, per cent: Korea and China 1960--2024, Japan
  1970--2024.}{World Bank, World Development Indicators (NY.GDS.TOTL.ZS)}
\begin{tikzpicture}
\begin{axis}[omrecord, width=0.84\linewidth, height=0.42\linewidth, xmin=1960, xmax=2024, ymin=0, ymax=55,
    xlabel={Year}, ylabel={Per cent of GDP}, xtick={1960,1970,1980,1990,2000,2010,2020},
    xticklabel style={/pgf/number format/1000 sep=}, unbounded coords=jump, legend style={at={(0.5,-0.3)},
    anchor=north, legend columns=3, /tikz/every even column/.append style={column sep=10pt}}]
  \addplot[omDef, thick] table[x=year, y=KOR] {figdata/bachelor-2/23-ramsey-model/out/saving.csv};
  \addplot[omThm, thick, dashed] table[x=year, y=JPN] {figdata/bachelor-2/23-ramsey-model/out/saving.csv};
  \addplot[omMeth, thick, dotted] table[x=year, y=CHN] {figdata/bachelor-2/23-ramsey-model/out/saving.csv};
  \legend{Korea, Japan, China}
\end{axis}
\end{tikzpicture}
\end{openingchart}

% figdata: saving_summary (KOR 1.0 in 1960, peak 41.7 in 1988, 34.0 in 2024; JPN 44.6 in 1970, 26.8 in 2024; CHN 37.0 in 1960, peak 50.6 in 2010, 43.3 in 2024)
Korea saved 1.0 per cent of its output in 1960 and 41.7 per cent in 1988; China's saving rate peaked at 50.6 per cent in 2010; Japan's fell from 44.6 per cent in 1970 to 26.8 in 2024.
A model in which households choose how much to save, looking ahead without limit, can produce saving rates that rise or
fall as an economy grows. It is the \omterm{def:b2:ramsey-model:ramsey}{Ramsey model}. Like the other models of this part, it is a toy: one household stands
for all, it lives for ever, and it knows the future. It is worth building because it makes saving a choice, gives the
benchmark against which the market's outcome can be judged, and shows how a dynamic economy responds to policy.

\begin{recall}
From the Year 1 volume: the Solow model, the steady state and the golden rule. From \cref{ch:b2:intertemporal-choice}: the \omterm{def:b2:intertemporal-choice:euler-equation}{Euler equation} and
the \omterm{def:b2:intertemporal-choice:eis}{elasticity of intertemporal substitution}. From the mathematics volumes: linear differential equations and the
eigenvalues of a $2 \times 2$ matrix.
\end{recall}

\begin{notation}[Symbols of this chapter]\label{not:b2:ramsey-model:symbols}
Time is continuous; $\dot x$ is the derivative of $x$ with respect to time. $k$ is capital per person, $c$ consumption per person, $f(k) = k^\alpha$ output per person;
$\rho$ the \omterm{def:b2:intertemporal-choice:time-preference}{rate of time preference}; $\theta$ the curvature of utility, $u(c) = c^{1-\theta}/(1 - \theta)$; $n$ population growth; $\delta$ depreciation; $r$ the real interest rate.
\end{notation}

\section{The household with an infinite horizon}\label{sec:b2:ramsey-model:household}

\begin{definition}[Representative household]\label{def:b2:ramsey-model:representative-household}
A \emph{representative household}\index{representative household} is a single household whose choices stand for the
total choices of all households in a model, as if all households were identical.
\end{definition}

\begin{definition}[Ramsey model]\label{def:b2:ramsey-model:ramsey}
The \emph{Ramsey model}\index{Ramsey model} (Ramsey--Cass--Koopmans) is a model of growth in which a \omterm{def:b2:ramsey-model:representative-household}{representative
household} with an infinite horizon chooses consumption and saving to maximise $\int_0^\infty e^{-(\rho - n)t}u(c(t))\,dt$, given the economy's production
function and capital accumulation.
\end{definition}

The household's \omterm{def:b2:intertemporal-choice:euler-equation}{Euler equation} of \cref{ch:b2:intertemporal-choice}, over a period of length $\Delta$ with \omterm{def:b2:intertemporal-choice:discount-factor}{discount factor} $e^{-\rho\Delta}$ and gross return $e^{r\Delta}$, reads $c(t)^{-\theta} = e^{-\rho\Delta}e^{r\Delta}c(t + \Delta)^{-\theta}$, so
$c(t + \Delta)/c(t) = e^{(r - \rho)\Delta/\theta}$. Taking logarithms, dividing by $\Delta$ and letting $\Delta \to 0$ gives the \omterm{def:b2:intertemporal-choice:euler-equation}{Euler equation} in continuous time:
\begin{equation}\label{eq:b2:ramsey-model:euler}
\frac{\dot c}{c} = \frac{r - \rho}{\theta}.
\end{equation}
Consumption grows when the interest rate exceeds the \omterm{def:b2:intertemporal-choice:time-preference}{rate of time preference}, and the more so the less curved utility is:
$1/\theta$ is the \omterm{def:b2:intertemporal-choice:eis}{elasticity of intertemporal substitution}. With $r = 0.05$, $\rho = 0.03$ and $\theta = 2$, consumption grows 1 per cent a year.

\begin{definition}[Transversality condition]\label{def:b2:ramsey-model:transversality}
The \emph{transversality condition}\index{transversality condition} requires that the present value of the household's
wealth tends to zero as time goes to infinity: an optimising household neither dies with unspent wealth in present value
nor runs up debt it never repays.
\end{definition}

The \omterm{def:b2:intertemporal-choice:euler-equation}{Euler equation} says how consumption changes; the budget constraint and the \omterm{def:b2:ramsey-model:transversality}{transversality condition} fix its level.
Without the condition, any starting consumption would satisfy the \omterm{def:b2:intertemporal-choice:euler-equation}{Euler equation}, and the household could save without
end or borrow without end.

\section{Capital accumulation}\label{sec:b2:ramsey-model:capital}

With firms that pay each factor its marginal product, the interest rate is $r = f'(k) - \delta$. Output not consumed is invested, and capital per
person grows by investment less depreciation and the dilution of population growth:
\begin{equation}\label{eq:b2:ramsey-model:k}
\dot k = f(k) - c - (n + \delta)k, \qquad \frac{\dot c}{c} = \frac{f'(k) - \delta - \rho}{\theta}.
\end{equation}
These two equations, one for capital and one for consumption, make up the model; we assume $\rho > n$.

\section{The phase diagram and the saddle path}\label{sec:b2:ramsey-model:phase}

\begin{definition}[Phase diagram]\label{def:b2:ramsey-model:phase-diagram}
A \emph{phase diagram}\index{phase diagram} of a system of two differential equations plots, in the plane of its two
variables, the curves on which each variable is constant and the direction in which the system moves in each region.
\end{definition}

The $\dot c = 0$ locus is the vertical line $f'(k^*) = \rho + \delta$; to its left the return is high and consumption rises. The $\dot k = 0$ locus is the curve $c = f(k) - (n + \delta)k$, highest at the
golden-rule capital stock; above it capital falls. Take $\alpha = 1/3$, $\delta = 0.05$, $\rho = 0.03$, $n = 0.01$ and $\theta = 2$. Then $k^* = (\alpha/(\rho + \delta))^{1/(1-\alpha)} = (0.333/0.08)^{1.5} = 8.505$, $c^* = 8.505^{1/3} - 0.06 \times 8.505 = 1.531$, and the saving rate is
$(n + \delta)k^*/f(k^*) = 0.25$.

\begin{definition}[Saddle path]\label{def:b2:ramsey-model:saddle-path}
The \emph{saddle path}\index{saddle path} of a system with a saddle-point steady state is the unique curve of initial
positions from which the system converges to the steady state; from every other position it diverges.
\end{definition}

\begin{omfigure}
\begin{tikzpicture}[econ]
\begin{axis}[omrecord, grid=none, xmin=0, xmax=20, ymin=0, ymax=2.6, width=0.7\linewidth, height=0.5\linewidth,
    xlabel={Capital per person, $k$}, ylabel={Consumption per person, $c$}, xtick={8.505,13.09},
    xticklabels={$k^*$,$k_{\mathrm{gold}}$}]
  \addplot[omThm, thick, domain=0.01:20, samples=100] {x^(1/3) - 0.06*x};
  \addplot[omDef, thick] coordinates {(8.505,0) (8.505,2.6)};
  \addplot[omMeth, very thick] table[x=k, y=c] {figdata/bachelor-2/23-ramsey-model/out/saddle.csv};
  \node[eqpoint] at (axis cs:8.505,1.531) {};
  \node[font=\footnotesize, anchor=west, omThm] at (axis cs:15.5,1.33) {$\dot k = 0$};
  \node[font=\footnotesize, anchor=west, omDef] at (axis cs:8.7,2.45) {$\dot c = 0$};
  \node[font=\footnotesize, anchor=west, omMeth] at (axis cs:15.5,2.2) {saddle path};
  \draw[->, thick] (axis cs:3,1.9) -- (axis cs:3,2.15);
  \draw[->, thick] (axis cs:3,1.9) -- (axis cs:2.2,1.9);
  \draw[->, thick] (axis cs:4,0.6) -- (axis cs:4,0.85);
  \draw[->, thick] (axis cs:4,0.6) -- (axis cs:4.8,0.6);
  \draw[->, thick] (axis cs:13,2.3) -- (axis cs:13,2.05);
  \draw[->, thick] (axis cs:13,2.3) -- (axis cs:12.2,2.3);
  \draw[->, thick] (axis cs:15,0.6) -- (axis cs:15,0.35);
  \draw[->, thick] (axis cs:15,0.6) -- (axis cs:15.8,0.6);
\end{axis}
\end{tikzpicture}

{\small The Ramsey \omterm{def:b2:ramsey-model:phase-diagram}{phase diagram} for the chapter's parameters ($\alpha = 1/3$, $\delta = 0.05$, $\rho = 0.03$, $n = 0.01$, $\theta = 2$). Above the $\dot k = 0$ curve capital falls; to the left of $\dot c = 0$ consumption rises. The \omterm{def:b2:ramsey-model:saddle-path}{saddle
path}, computed numerically, is the only route to the steady state; it lies to the left of the golden-rule stock
$k_{\mathrm{gold}} = 13.09$.}
\end{omfigure}

\begin{proposition}[The steady state is a saddle point]\label{prop:b2:ramsey-model:saddle}
Linearised around $(k^*, c^*)$, the system \eqref{eq:b2:ramsey-model:k} has one positive and one negative eigenvalue.
\end{proposition}

\begin{proof}
The Jacobian at the steady state has entries $\partial\dot k/\partial k = f'(k^*) - (n + \delta) = \rho - n$, $\partial\dot k/\partial c = -1$, $\partial\dot c/\partial k = c^*f''(k^*)/\theta$ and $\partial\dot c/\partial c = (f'(k^*) - \delta - \rho)/\theta = 0$. Its determinant is
$0 - (-1)c^*f''(k^*)/\theta = c^*f''(k^*)/\theta < 0$, since $f'' < 0$. The determinant is the product of the eigenvalues, so they are real and of opposite signs.
\end{proof}

With the chapter's numbers the Jacobian is $\begin{pmatrix} 0.02 & -1 \\ -0.0048 & 0 \end{pmatrix}$, with eigenvalues 0.08 and $-0.06$. Near the steady state the economy closes 6 per cent of its gap a year, so
half the gap in $\ln 2/0.06 = 11.6$ years. Capital cannot jump; consumption can. Given $k(0)$, the household picks the one $c(0)$ on the \omterm{def:b2:ramsey-model:saddle-path}{saddle path}; any other choice
leads either to capital running out or to capital piling up beyond the golden rule, which violates the \omterm{def:b2:ramsey-model:transversality}{transversality
condition}.

\omcode{code/bachelor-2/23-ramsey-model/listings/ramsey.py}{1}{37}{Shooting for the \omterm{def:b2:ramsey-model:saddle-path}{saddle path} from a quarter of the
steady-state capital, for two curvatures of utility.}\label{lst:b2:ramsey-model:shooting}

\omoutput{code/bachelor-2/23-ramsey-model/outputs/ramsey.txt}{With $\theta = 2$ the saving rate falls along the way; with $\theta = 6$ it rises; the
convergence is slower with the larger curvature.}\label{out:b2:ramsey-model:shooting}

% code: listings/ramsey.py
The listing finds the \omterm{def:b2:ramsey-model:saddle-path}{saddle path} by trying initial consumption levels and keeping the one that neither exhausts capital
nor overshoots. The model makes the saving rate along the path depend on the curvature of utility: with $\theta = 2$ a poor economy saves 34 per cent at first and less
as it grows, since households that substitute easily over time take advantage of the high return; with $\theta = 6$, households
reluctant to accept low consumption today start at 21 per cent and save more as they grow richer.

\begin{omfigure}
\begin{tikzpicture}
\begin{axis}[omrecord, name=kpath, width=0.46\linewidth, height=0.4\linewidth, xmin=0, xmax=60, ymin=0, ymax=1.05,
    xlabel={Years}, ylabel={$k/k^*$}, legend style={at={(0.95,0.05)}, anchor=south east, font=\footnotesize}]
  \addplot[omDef, thick] table[x=t, y=k2] {figdata/bachelor-2/23-ramsey-model/out/paths.csv};
  \addplot[omThm, thick, dashed] table[x=t, y=k6] {figdata/bachelor-2/23-ramsey-model/out/paths.csv};
  \legend{$\theta = 2$, $\theta = 6$}
\end{axis}
\begin{axis}[omrecord, at={(kpath.east)}, anchor=west, xshift=1.6cm, width=0.46\linewidth, height=0.4\linewidth,
    xmin=0, xmax=60, ymin=0.15, ymax=0.4, xlabel={Years}, ylabel={Saving rate},
    yticklabel style={/pgf/number format/fixed, /pgf/number format/precision=2}]
  \addplot[omDef, thick] table[x=t, y=s2] {figdata/bachelor-2/23-ramsey-model/out/paths.csv};
  \addplot[omThm, thick, dashed] table[x=t, y=s6] {figdata/bachelor-2/23-ramsey-model/out/paths.csv};
\end{axis}
\end{tikzpicture}

{\small Model simulation, the chapter's parameters, starting from a quarter of the steady-state capital: capital relative to
its steady state (left) and the saving rate (right), for $\theta = 2$ and $\theta = 6$. Invented parameters, not data.}
\end{omfigure}

\begin{model}{The Ramsey model}
\assumes A \omterm{def:b2:ramsey-model:representative-household}{representative household} with an infinite horizon and perfect foresight, competitive firms with diminishing
returns to capital, and no frictions in adjusting capital.
\predicts Convergence along the \omterm{def:b2:ramsey-model:saddle-path}{saddle path} to a steady state set by the \omterm{def:b2:ramsey-model:modified-golden-rule}{modified golden rule}; a saving rate that falls or
rises along the way depending on the curvature of utility; no response of steady-state capital to permanent government
purchases.
\tested Poor economies with access to the same technology grow faster; but the convergence the model implies, about 6
per cent of the gap a year here, is faster than the data on American states show, which can be matched only if
diminishing returns set in very slowly (Barro and Sala-i-Martin 1992).
\breaks One household cannot represent households that differ in wealth and age; without \omterm{def:b2:investment:adjustment-cost}{adjustment costs} or slow
diffusion of technology, convergence is too fast; the large swings in saving rates of the opening chart need more than
the model's curvature of utility.
\end{model}

\section{The modified golden rule}\label{sec:b2:ramsey-model:golden}

\begin{definition}[Modified golden rule]\label{def:b2:ramsey-model:modified-golden-rule}
The \emph{modified golden rule}\index{modified golden rule} is the condition $f'(k^*) = \rho + \delta$ that fixes steady-state capital in the \omterm{def:b2:ramsey-model:ramsey}{Ramsey model}: the
marginal product of capital, net of depreciation, equals the \omterm{def:b2:intertemporal-choice:time-preference}{rate of time preference}.
\end{definition}

The golden rule of the Solow model, $f'(k) = n + \delta$, maximises steady-state consumption. Because $\rho > n$, the \omterm{def:b2:ramsey-model:modified-golden-rule}{modified golden rule} stops short of it: here $k^* = 8.505$ against $k_{\mathrm{gold}} = 13.09$. Reaching the golden rule
would require consumption to fall today to build capital whose payoff comes later, and an impatient household does not
think the trade worth it. The economy is therefore never \omterm{def:b2:overlapping-generations:dynamic-inefficiency}{dynamically inefficient}, unlike the overlapping-generations
economy of \cref{ch:b2:overlapping-generations}: an infinitely lived household would never hold capital that earns less than the growth rate.

\section{The market equals the planner}\label{sec:b2:ramsey-model:planner}

\begin{proposition}[The competitive path is the planner's path]\label{prop:b2:ramsey-model:planner}
A planner who maximises the \omterm{def:b2:ramsey-model:representative-household}{representative household}'s utility subject only to $\dot k = f(k) - c - (n + \delta)k$ chooses the same path of $k$ and $c$ as households
and firms in competitive markets.
\end{proposition}

\begin{proof}
The planner's optimal consumption satisfies an \omterm{def:b2:intertemporal-choice:euler-equation}{Euler equation} in which the return on saving is the net marginal product
of capital: $\dot c/c = (f'(k) - \delta - \rho)/\theta$. In the market, the household's \omterm{def:b2:intertemporal-choice:euler-equation}{Euler equation} is $\dot c/c = (r - \rho)/\theta$ and firms set $r = f'(k) - \delta$, which gives the same equation. The
household's budget, wage plus interest income less consumption, is $\dot k = w + rk - c - nk$, and with $w = f(k) - kf'(k)$ and $r = f'(k) - \delta$ it becomes $\dot k = f(k) - c - (n + \delta)k$, the planner's constraint. Both
paths satisfy the same two equations and the same \omterm{def:b2:ramsey-model:transversality}{transversality condition}, from the same $k(0)$, so they coincide.
\end{proof}

The proposition is the first welfare theorem of \cref{ch:b2:welfare-theorems} in a dynamic economy. It holds because the model has no
\omterm{def:b2:externalities-and-public-goods-theory:externality}{externality}, no market power and one household; each of these, removed, gives a reason the market's path can differ from
the planner's.

\section{A change in government purchases}\label{sec:b2:ramsey-model:government}

Let the government buy $G$ per person, financed by lump-sum taxes: $\dot k = f(k) - c - G - (n + \delta)k$. The $\dot k = 0$ curve shifts down by $G$; the $\dot c = 0$ line does not move, since the \omterm{def:b2:ramsey-model:modified-golden-rule}{modified golden rule}
does not involve $G$.

\begin{omfigure}
\begin{tikzpicture}[econ]
\begin{axis}[omrecord, grid=none, xmin=0, xmax=20, ymin=0, ymax=2.6, width=0.62\linewidth, height=0.46\linewidth,
    xlabel={$k$}, ylabel={$c$}, xtick={8.505}, xticklabels={$k^*$}]
  \addplot[omThm, thick, domain=0.01:20, samples=100] {x^(1/3) - 0.06*x};
  \addplot[omThm, thick, dashed, domain=0.3:20, samples=100] {x^(1/3) - 0.06*x - 0.3};
  \addplot[omDef, thick] coordinates {(8.505,0) (8.505,2.6)};
  \node[eqpoint] at (axis cs:8.505,1.531) {};
  \node[eqpoint] at (axis cs:8.505,1.231) {};
  \node[font=\footnotesize, anchor=west] at (axis cs:8.7,1.62) {$E_0$};
  \node[font=\footnotesize, anchor=west] at (axis cs:8.7,1.14) {$E_1$};
  \draw[->, thick, black!70] (axis cs:8.2,1.5) -- (axis cs:8.2,1.27);
  \node[font=\footnotesize, anchor=west, omThm] at (axis cs:15.5,1.36) {$G = 0$};
  \node[font=\footnotesize, anchor=west, omThm] at (axis cs:15.5,0.92) {$G > 0$};
\end{axis}
\end{tikzpicture}

{\small A permanent rise in government purchases. The $\dot k = 0$ curve shifts down by $G$; consumption drops at once by $G$, from $E_0$ to $E_1$, and capital does not move:
the household pays for the purchases entirely out of consumption.}
\end{omfigure}

A permanent rise in $G$ lowers consumption one for one at once, and nothing else changes: the economy jumps from one steady state
to the other. A temporary rise is different: households, knowing it will end, cut consumption by less than $G$ and let capital fall
for a while, then rebuild it; the interest rate rises during the episode. That is the model's version of crowding out.

\begin{record}[Korea's saving and capital deepening]
% figdata: ky_summary (1980 3.09, min 2.98 in 1988, 2023 5.11); initial_share (1960 0.847, 1970 0.410, 1980 0.104); saving_summary
Korea's saving rate rose from 1 per cent of GDP in 1960 to 41.7 per cent in 1988 and has stayed above 30 per cent since. Its capital
stock was 3.09 times a year's output in 1980 and, after a flat decade, rose faster than output, to 5.11 times in 2023: the capital deepening that the transition of the \omterm{def:b2:ramsey-model:ramsey}{Ramsey model} describes, though over a much longer
time than the model's half-life of twelve years.
\end{record}

\begin{omchart}{Capital stock relative to GDP, Korea, 1980--2023, constant national prices. The series starts from a
  capital stock assumed for 1953, which was still 85 per cent of the 1960 stock and 41 per cent of the 1970 one, so the
  earlier years show that guess wearing out rather than data; by 1980 it is a tenth of the stock.}{Penn World Table
  (rnna, rgdpna, delta); our computation}
\begin{tikzpicture}
\begin{axis}[omrecord, width=0.72\linewidth, height=0.4\linewidth, xmin=1980, xmax=2023, ymin=0, ymax=6,
    xlabel={Year}, ylabel={Capital / GDP}, xtick={1980,1990,2000,2010,2020},
    xticklabel style={/pgf/number format/1000 sep=}]
  \addplot[omDef, thick] table[x=year, y=ky] {figdata/bachelor-2/23-ramsey-model/out/ky.csv};
\end{axis}
\end{tikzpicture}
\end{omchart}

\section{Exercises}\label{sec:b2:ramsey-model:exercises}

\begin{exercise}[$\star$]\label{exo:b2:ramsey-model:1}
Read the opening chart. Describe Korea's saving rate from 1960 to 2024, and compare it with Japan's.
\end{exercise}

\begin{exercise}[$\star$]\label{exo:b2:ramsey-model:2}
With $\alpha = 1/3$, $\delta = 0.05$, $\rho = 0.03$ and $n = 0.01$, compute $k^*$, $y^*$, $c^*$ and the saving rate.
\end{exercise}

\begin{exercise}[$\star$]\label{exo:b2:ramsey-model:3}
Read the chart of Korea's capital--output ratio. When was it lowest, and what has happened since? Why does the chart
not start in 1960?
\end{exercise}

\begin{exercise}[$\star$]\label{exo:b2:ramsey-model:4}
What evidence on the growth of poor and rich regions would contradict the model's prediction of fast convergence?
\end{exercise}

\begin{exercise}[$\star\star$]\label{exo:b2:ramsey-model:5}
Compute the golden-rule capital stock for the parameters of Exercise 2 and compare it with $k^*$. Why is the steady state below
it?
\end{exercise}

\begin{exercise}[$\star\star$]\label{exo:b2:ramsey-model:6}
Compute the Jacobian at the steady state, its eigenvalues and the half-life of the gap, for $\theta = 2$ and for $\theta = 6$.
\end{exercise}

\begin{exercise}[$\star\star$]\label{exo:b2:ramsey-model:7}
Predict, then compute: households become more impatient, $\rho = 0.04$. What happens to $k^*$, $c^*$, the saving rate, and to consumption at the moment
of the change?
\end{exercise}

\begin{exercise}[$\star\star$]\label{exo:b2:ramsey-model:8}
Use the \omterm{def:b2:ramsey-model:phase-diagram}{phase diagram} to describe the response to a rise in $G$ that is announced today and takes effect, permanently, in five
years.
\end{exercise}

\begin{exercise}[$\star\star$]\label{exo:b2:ramsey-model:9}
With $r = 0.05$, $\rho = 0.03$ and $\theta = 2$, at what rate does consumption grow? And with $\theta = 1$?
\end{exercise}

\begin{exercise}[$\star\star\star$]\label{exo:b2:ramsey-model:10}
Derive \eqref{eq:b2:ramsey-model:euler} from the discrete \omterm{def:b2:intertemporal-choice:euler-equation}{Euler equation} $u'(c_t) = \beta(1 + r)u'(c_{t+1})$ with $u(c) = c^{1-\theta}/(1 - \theta)$, taking periods of length $\Delta$ and letting $\Delta \to 0$.
\end{exercise}

\begin{exercise}[$\star\star\star$]\label{exo:b2:ramsey-model:11}
Show that the eigenvalues of $\begin{pmatrix} a & -1 \\ b & 0 \end{pmatrix}$ are $\tfrac12(a \pm \sqrt{a^2 - 4b})$, and that with $b < 0$ the speed of convergence rises with $|b|$. What lowers $|b|$?
\end{exercise}

\begin{exercise}[$\star\star\star$]\label{exo:b2:ramsey-model:12}
Suppose capital income is taxed at rate $\tau$, the revenue returned lump-sum. Show that the steady state satisfies $(1 - \tau)(f'(k) - \delta) = \rho$, and that $k^*$ falls with $\tau$.
\end{exercise}

\section{Problem: Kestria's catch-up}\label{sec:b2:ramsey-model:problem}

\begin{problem}[{Weekend problem --- Kestria's catch-up: the household, the steady state, the speed of convergence, and a
permanent rise in government purchases}]
\label{pb:b2:ramsey-model:1}
Kestria's economy is a Ramsey economy with $f(k) = k^{0.3}$, $\delta = 0.06$, $\rho = 0.04$, $n = 0$ and $\theta = 2$; output and consumption are in thousands of kests per person a
year.

\medskip
\noindent\textbf{Part I --- The household.}
\begin{enumerate}
  \item Write the household's \omterm{def:b2:intertemporal-choice:euler-equation}{Euler equation}.
  \item If the interest rate is 6 per cent, how fast does consumption grow?
  \item What is the \omterm{def:b2:intertemporal-choice:eis}{elasticity of intertemporal substitution}?
  \item Why does the household need the \omterm{def:b2:ramsey-model:transversality}{transversality condition}?
  \item Write the law of motion of capital.
\end{enumerate}

\medskip
\noindent\textbf{Part II --- The steady state.}
\begin{enumerate}[resume]
  \item What interest rate holds in the steady state?
  \item Find $k^*$.
  \item Find output, consumption and the wage.
  \item Find the saving rate.
  \item Find the golden-rule capital stock and compare.
  \item Is Kestria saving too little?
\end{enumerate}

\medskip
\noindent\textbf{Part III --- Convergence.}
\begin{enumerate}[resume]
  \item Write the Jacobian at the steady state.
  \item Find its eigenvalues.
  \item What is the half-life of a gap in capital?
  \item With $\theta = 4$, what is the half-life?
  \item Why does a larger $\theta$ slow convergence?
  \item Kestria starts with a quarter of $k^*$. Does its saving rate start above or below 0.18 when $\theta = 2$?
  \item Compare the model's half-life with the evidence of the model box.
\end{enumerate}

\medskip
\noindent\textbf{Part IV --- Government purchases.} Kestria's government permanently raises purchases by 0.1 per person,
financed by lump-sum taxes.
\begin{enumerate}[resume]
  \item What happens to $k^*$?
  \item What happens to consumption, and when?
  \item What happens to the interest rate?
  \item If the rise were for five years only, would consumption fall by more or by less?
  \item If the purchases were financed by a tax on capital income, would $k^*$ change?
  \item What does the model leave out that could make government purchases raise output in the short run?
  \item State the result: steady-state capital, the saving rate, and the half-life of convergence.
\end{enumerate}
\end{problem}
